% !TEX root = ../main.tex

\chapter{Section stacks and derived solution stacks}
\label[chapter]{chap:Section_Stacks_Jets_Solution_Stacks}

A differential operator assigns to a function, section, or map an expression
in its derivatives. Its solutions often describe a geometric condition:
harmonicity, stationarity of an energy, or compatibility with a complex
structure. The equation can be defined for smooth families of unknowns.
The resulting functor of solutions is a stack, whose geometric
representability is a separate question.

Our running example is the nonlinear elliptic equation
\[
 P(u)=\Delta u+u^2=0
\]
on a closed connected Riemannian manifold. Its ordinary solution set is
a single point, but its infinitesimal solutions and obstructions are
both nontrivial. The eventual conclusion is concrete: near zero, its
derived solution object is the multiple point defined by $c^2=0$.
\Cref{chap:Elliptic_Representability_Derived_Charts} will prove this by
eliminating the mean-zero part of $u$, leaving its constant part $c$
as the remaining variable.

This chapter sets up the functor that the later analytic argument will
represent. Geometric examples lead to section stacks and their derived
fibres. Jets then express the finite-order condition, and linearization
gives the definition of ellipticity. Each construction is formulated
for families, so it can be used over a parameter stack.

\section{Differential equations from geometry}

The examples below identify the unknowns, their equations, and the
bundles in which the equations take values. We begin with scalar
functions and then consider maps, where the output bundle depends
on the unknown.

\begin{example}[Harmonic functions]
For a Riemannian manifold $N$, the Laplacian is an operator
\[
 \Delta\colon\Cinfty(N)\longrightarrow\Cinfty(N).
\]
We use the nonnegative convention $\Delta=d^*d$.
The equation $\Delta u=0$ says that $u$ is harmonic.
It occurs in potential theory and in extending boundary data into a domain.
On a closed connected manifold, integration by parts shows that harmonic
functions are constant:
\[
 0=\int_Nu\Delta u=\int_N|du|^2.
\]
Boundary-value problems require additional analytic hypotheses; the
representability theorem considered here concerns closed fibres.

Both source and target are sections of the trivial real line bundle
$N\times\R\to N$. The operator takes an entire section as its input
and returns another section.
\end{example}

\begin{example}[Critical points of energies]
Let $N$ be a closed Riemannian manifold and let $V\colon\R\to\R$ be smooth. Consider
\[
 \mathcal E(u)=\int_N\left(\tfrac12|du|^2+V(u)\right)\,d\mathrm{vol}.
\]
Differentiating in the direction $v$ and integrating by parts gives
\[
 d\mathcal E_u(v)=\int_N\bigl(\Delta u+V'(u)\bigr)v\,d\mathrm{vol}.
\]
Therefore the critical points are the solutions of
$P(u)=0$, where
\[
 P\colon\Cinfty(N)\longrightarrow\Cinfty(N),\qquad
 P(u)=\Delta u+V'(u).
\]
The energy is a real-valued functional; its Euler--Lagrange operator $P$
has a function as its output. The metric and volume form identify the
first variation with this function.

Taking $V(u)=u^3/3$ gives the running example $P(u)=\Delta u+u^2$.
\end{example}

\begin{example}[Pseudo-holomorphic maps]
Let $(\Sigma,j)$ be a fixed Riemann surface and $(Z,J)$ an almost complex manifold.
A smooth map $u\colon\Sigma\to Z$ is pseudo-holomorphic when
\[
 \bar\partial_Ju=\tfrac12(du+J(u)\circ du\circ j)=0.
\]
The unknown is a section of $\Sigma\times Z\to\Sigma$.
Its output lies in
\[
 \Gamma\bigl(\Sigma,
 \operatorname{Hom}^{0,1}_{\R}(T\Sigma,u^*TZ)\bigr),
\]
where the superscript denotes complex-antilinear real linear maps.
The output bundle depends on $u$.

Globally, this is a section of a bundle over the space of maps.
Near a fixed map, exponential coordinates on $Z$ and identifications
of nearby output bundles turn it into an operator
\[
 Q\subseteq\Gamma(\Sigma,F)\longrightarrow\Gamma(\Sigma,E)
\]
for fixed finite-rank bundles $F,E$.
\Cref{chap:Elliptic_Equations_Kuranishi_Models} explains the coordinate principle, and \Cref{chap:Pseudo_Holomorphic_Curves_Derived_Moduli} carries out
the Cauchy--Riemann application. Keeping track of these source and target
bundles is necessary before applying an ellipticity theorem.
\end{example}

\section{Section stacks}

To apply the functor-of-points viewpoint, we retain smooth families of
sections and their pullbacks. We first fix the domain manifold, then
allow it to vary over a base.

For a submersion $p\colon X\to N$, the ordinary section set is
\[
 \Gamma(N;X)=\{u\colon N\to X:\ p\circ u=\operatorname{id}_N\}.
\]
If $X$ is a vector bundle, this is a vector space.
A family parametrized by an ordinary manifold $T$ is a smooth map
$T\times N\to X$ over $N$. Smoothness is joint smoothness in the
parameter and in the point of $N$.

The same formula makes sense with derived test objects.
For $T$ derived, it uses maps of derived stacks, and so retains
infinitesimal and higher data in the parameter.

\begin{definition}[Section stack]
\label[definition]{def:Section_Stack}
The section stack of $X\to N$ is
\[
 \Sec_N(X)(T)=\Hom_{/N}(T\times N,X).
\]
\end{definition}

This is the mapping stack often written $\operatorname{Map}_N(N,X)$.
The notation $\Gamma(N;X)$ will be reserved for ordinary smooth
sections, and, for a vector bundle on compact $N$, their Fr\'echet
manifold. The section stack is a functor on parameter objects.

A family of manifolds over a smooth stack $S$
is a map $q\colon M\to S$ whose pullback along every ordinary manifold
$T\to S$ is a submersion of ordinary manifolds. The family is proper
if each such pullback is proper. For a manifold base this is the usual
notion of a submersion, or a proper submersion, respectively.
Suppose $q\colon M\to S$ is such a family and $X\to M$ is a family
of submersions.

\begin{definition}[Relative section stack]
\label[definition]{def:Relative_Section_Stack}
The relative section stack $\Sec_{M/S}(X)\to S$ is characterized by
\[
 \Hom_{/S}(T,\Sec_{M/S}(X))
 \cong\Hom_{/M}(T\times_S M,X).
\]
It is the Weil restriction $\operatorname{Res}_{M/S}(X)$.
\end{definition}

The terminology describes the right adjoint to pullback along $M\to S$.
Its existence here follows from the locally cartesian closed
$\infty$-category of stacks
\cite[Definition~3.2.1]{Steffens_Representability_Elliptic_Moduli}.
The formula also proves base-change compatibility:
\[
 \Sec_{M/S}(X)\times_S T
 \cong\Sec_{M_T/T}(X_T),
 \qquad M_T=M\times_S T.
\]
Over $s\in S$, its fibre is the section stack on $M_s$.

\section{Solution stacks as derived fibres}

Let $P\colon\Sec_N(X)\to\Sec_N(Y)$ be any morphism of stacks, and let
$b$ be a section of $Y$. Define
\begin{equation}\label{eq:Solution_Stack_Inhomogeneous}
 \Sol(P;b)=\Sec_N(X)\times_{\Sec_N(Y)}^{\mathrm{der}}*,
\end{equation}
where the second map classifies $b$.
The adjective derived emphasizes the category in which the pullback is
taken; all pullbacks in derived stacks are already homotopy pullbacks.

On a derived parameter $T$, a solution is a family $u$ together with a
homotopy from $P(u)$ to the constant family $b_T$.
Thus a solution stack can be defined without any finite-order assumption.
One does need a morphism of stacks, not merely a map between sets
of smooth sections.

For closed $N$, a smooth map between the corresponding convenient
section manifolds supplies such a morphism. This follows from the
comparison between section stacks and convenient manifolds explained
in \Cref{chap:Elliptic_Representability_Smooth_Derived_Bridge}. For example, on $S^1$ the map
\[
 u\longmapsto\left(x\longmapsto\int_{S^1}u\,d\theta\right)
\]
is a smooth linear map of section manifolds and defines a solution stack.
It is not a finite-order differential operator: changing $u$ away from
$x$ can change the output at $x$ while leaving every derivative there fixed.

For a parameter base $S$, a right-hand side is a section
$b\colon S\to\Sec_{M/S}(Y)$, and the same pullback is formed over $S$.
All these constructions commute with base change.
The finite-order differential condition introduced next provides a
geometric way to define $P$ and the analytic structure needed for
representability.

\section{Jets and finite-order operators}

Jets encode the dependence of an operator on finitely many derivatives.
Their relative construction also makes the induced map of section
stacks available on derived parameters.

A $k$-jet of a local section of $X\to N$ at $x\in N$ records its value
and its derivatives through order $k$ at $x$.
Two sections define the same jet when those derivatives agree in local
coordinates. The chain rule makes this independent of the coordinates.
These jets form the bundle $J_N^k(X)\to N$.

For example, for $X=N\times\R$, a first jet records a scalar and a covector.
A second jet records the second derivatives as well; a connection
identifies that last datum with a symmetric tensor.
The jet prolongation sends a section $u$ to $j^ku$.

\begin{definition}[Finite-order differential operator]
\label[definition]{def:Finite_Order_Differential_Operator}
A nonlinear differential operator of order at most $k$ from $X$ to $Y$
over $N$ is a smooth map
\[
 p\colon J_N^k(X)\longrightarrow Y
\]
over $N$. Its induced map on sections is
$P(u)=p\circ j^ku$.
\end{definition}

Using the lowercase $p$ distinguishes the finite-dimensional map on
jets from the induced map $P$ on sections.
The latter is determined at $x$ by finitely many derivatives of $u$ at
that same point. A general map on sections need not have this locality.

For $\Delta u+u^2$, $k=2$. In local coordinates the Laplacian depends
on derivatives of $u$ of order at most two, while $u^2$ uses only its value.
For the Cauchy--Riemann equation only first jets are needed.

For $M\to S$, relative jets $J^k_{M/S}(X)$ record derivatives in the
directions tangent to the fibres $M_s$. No differentiation in the
parameter $s$ is involved. The resulting map $P$ is therefore an
$S$-family of operators on the fibre manifolds.

The derived jet stack has a coordinate-independent construction.
Let $(M/S)^{(k)}$ be the order-$k$ infinitesimal neighbourhood of the
relative diagonal, with projections $\pi_1,\pi_2$ to $M$. Then
\[
 J^k_{M/S}(X)
 =\operatorname{Res}_{(M/S)^{(k)}/M}
       \bigl(X\times_{M,\pi_1}(M/S)^{(k)}\bigr),
\]
where the Weil restriction uses $\pi_2$.
It parametrizes sections on infinitesimal neighbourhoods.
For an ordinary submersion this recovers the classical jet bundle.
The construction, the prolongation map, and base-change compatibility are
proved in
\cite[Definition~3.3.7 and Proposition~3.3.9]
{Steffens_Representability_Elliptic_Moduli}.

In particular, a map $p\colon J^k_{M/S}(X)\to Y$ induces a morphism
of relative section stacks. This makes the preceding solution
construction available over all derived test objects.

\section{Differential moduli problems and ellipticity}

We now collect the geometric data of an equation into a differential
moduli problem. Its linearization gives the deformation-obstruction
complex, and its principal symbol determines whether elliptic analysis
can be applied.

A useful common formulation asks two operators to have the same output:
\[
 P_1(u_1)=P_2(u_2).
\]
It includes $P(u)=b$ by taking the second unknown bundle to be $N\to N$,
whose only section is the identity, and taking its operator to be $b$.

\begin{definition}
\label[definition]{def:Differential_Moduli_Problem}
Over a family $M\to S$, a differential moduli problem consists of
submersions $Y_1,Y_2,X\to M$ and maps
\[
 p_i\colon J^k_{M/S}(Y_i)\longrightarrow X.
\]
Its solution stack is
\[
 \Sol(\mathcal E)=
 \Sec_{M/S}(Y_1)
 \times_{\Sec_{M/S}(X)}
 \Sec_{M/S}(Y_2),
\]
with maps induced by $p_1,p_2$.
\end{definition}

This is \cite[Definitions~3.4.1--3.4.2]
{Steffens_Representability_Elliptic_Moduli}.
The order $k$ is a common order: an operator of lower order can be
composed with a jet truncation map. Ellipticity below is then a
condition on the symbol for that chosen common order.

When $S$ is a smooth stack, these data are specified on manifold
families over $S$ and are compatible with pullback. Relative jets and
section stacks commute with that pullback, so the equations and
solution stacks descend. This includes families with symmetries
encoded in the base stack.

Fix first $S=*$ and a solution $u$ of $P(u)=b$.
Variations of $u$ are sections of
$F=u^*T_{X/N}$; variations of the output at $b$ are sections of
$E=b^*T_{Y/N}$. Differentiating the finite jet formula gives a linear
differential operator
\[
 D_uP\colon\Gamma(N;F)\longrightarrow\Gamma(N;E).
\]
In coordinates it has the form
\[
 D_uP(v)=\sum_{|\alpha|\leq k}A_\alpha(x)\partial^\alpha v.
\]
Its coefficients depend on $u$ and its derivatives, but the expression
is linear in $v$.

For a matching problem and solution $(u_1,u_2)$ with common output $w$,
put
\[
 F=u_1^*T_{Y_1/N}\oplus u_2^*T_{Y_2/N},\qquad
 E=w^*T_{X/N}.
\]
The linearized matching equation is
\begin{equation}\label{eq:Tangent_Complex_Solution_Stack}
 D(v_1,v_2)=D_{u_1}P_1(v_1)-D_{u_2}P_2(v_2).
\end{equation}
Its deformation-obstruction complex is
$[\Gamma(N;F)\xrightarrow D\Gamma(N;E)]$ in degrees $0,-1$.
The kernel is the infinitesimal solution space.
The cokernel receives obstructions to solving inhomogeneous correction
equations, as in the finite-dimensional zero-locus calculation.
The representability proof will identify this complex with the
finite-dimensional tangent complex of a local derived chart.

Ellipticity is a condition on the highest-order part of the linearized
operator. In particular, it depends on the source and target bundles
of the full matching equation.

\begin{definition}[Principal symbol and ellipticity]
\label[definition]{def:Elliptic_Differential_Moduli_Problem}
For a linear operator $D$ of order $k$, its principal symbol is
\[
 \sigma_D(x,\xi)=\sum_{|\alpha|=k}A_\alpha(x)\xi^\alpha
 \colon F_x\longrightarrow E_x.
\]
It is elliptic if this map is invertible for every nonzero covector $\xi$.
A differential moduli problem is elliptic if its linearized matching
operator is elliptic at every solution on every fibre $M_s$.
\end{definition}

This real symbol convention omits the optional factors of $i$ used in
Fourier analysis; the invertibility condition is unchanged.
For the nonnegative Laplacian the real symbol with this convention is
$-|\xi|^2$, while the Fourier convention gives $|\xi|^2$.
Both are invertible away from the zero covector.

Ellipticity concerns the highest-order part. On a compact manifold
it leads to finite-dimensional kernel and cokernel after Sobolev
completion. It does not assert surjectivity. The latter failure is
exactly why a derived zero locus can be needed.

\section{Solutions and linearization of the running example}

Return to $P(u)=\Delta u+u^2$ on a closed connected manifold $N$.
If $P(u)=0$, integration gives
\[
 0=\int_N(\Delta u+u^2)=\int_Nu^2,
\]
so $u=0$. Its ordinary solution set is one point.

Differentiating gives
\[
 D_uP(v)=\Delta v+2uv.
\]
The zeroth-order term does not change the principal symbol, so this
operator is elliptic. At zero it is $\Delta$.
The standard solvability theorem for the Laplacian on a closed connected
manifold says that its image consists of the functions of integral zero.
Consequently,
\[
 \ker\Delta\cong\R,\qquad\operatorname{coker}\Delta\cong\R.
\]
Constants give representatives in both cases.

This explains the usefulness of the example. Its real zero set is
particularly simple, while its infinitesimal solutions and obstructions
are both nontrivial. \Cref{chap:Elliptic_Representability_Derived_Charts} will obtain its local Kuranishi equation
explicitly, recovering a multiple point.

\section{Further reading}

The jet tower makes the highest-order part of an operator explicit.
We describe it for vector bundles, then illustrate parameter dependence
with a family of metrics and give two further calculations.

For a vector bundle $F\to N$, the jet tower has exact sequences
\[
 0\longrightarrow\operatorname{Sym}^k(T^*N)\otimes F
 \longrightarrow J_N^k(F)\longrightarrow J_N^{k-1}(F)
 \longrightarrow0.
\]
The kernel consists of the highest derivatives when all lower derivatives
vanish. Connections can split these sequences, but the exact sequences
themselves are canonical. This is also the part of the jet data
used by the principal symbol.

A smooth family of metrics $g_s$ gives operators
$P_s(u)=\Delta_{g_s}u+u^2$ on the fibres. For a test family
$T\to S$, the equation is evaluated using the pulled-back metric.
The parameter affects the coefficients, while all derivatives of $u$
remain in the fibre directions.
The relative solution stack is defined before any assertion about
its projection to $S$ being a submersion or being representable.

\begin{exercise}
Compute the Euler--Lagrange operator for
$\int_N(\tfrac12|du|^2+\tfrac14u^4-\tfrac12u^2)$ and its linearization
at each constant solution. Check ellipticity.
\end{exercise}

\begin{exercise}
Let $P_1,P_2$ be first-order linear operators with a common target.
Write the principal symbol of the matching equation.
Explain why ellipticity of either operator individually does not
imply ellipticity of their difference on the direct sum of the sources.
\end{exercise}

The section and jet constructions are in
\cite[Sections~3.2--3.3 and Definitions~3.4.1--3.4.2]
{Steffens_Representability_Elliptic_Moduli}.
The remainder of the seminar proves the representability statement
for the elliptic problems defined here.
